\documentclass{article}
\usepackage[T1]{fontenc}
\usepackage{lmodern}
\usepackage{graphicx}
\usepackage{amsmath,amssymb}
\usepackage[a4paper,margin=2cm]{geometry}
\usepackage[absolute,overlay]{textpos}
\setlength{\TPHorizModule}{1mm}
\setlength{\TPVertModule}{1mm}
\usepackage[indonesian]{babel}
\usepackage{hyperref}
\setlength{\emergencystretch}{2em}
% unit: Halaman 23

\begin{document}

% === Halaman 23 (python/native) ===

S5: 

2 

12 

4 

1 

7 

10 

11 

6 

8 

5 

3 

15 

13 

0 

14 

9 

14 

11 

2 

12 

4 

7 

13 

1 

5 

0 

15 

10 

3 

9 

8 

16 

4 

2 

1 

11 

10 

13 

7 

8 

15 

9 

12 

5 

6 

3 

0 

14 

11 

8 

12 

7 

1 

14 

2 

13 

6 

15 

0 

9 

10 

4 

5 

3 



S6: 

12 

1 

10 

15 

9 

2 

6 

8 

0 

13 

3 

4 

14 

7 

5 

11 

10 

15 

4 

2 

7 

12 

9 

5 

6 

1 

13 

14 

0 

11 

3 

8 

9 

14 

15 

5 

2 

8 

12 

3 

7 

0 

4 

10 

1 

13 

11 

6 

4 

3 

2 

12 

9 

5 

15 

10 

11 

14 

1 

7 

6 

0 

8 

13 


S7: 

4 

11 

2 

14 

15 

0 

8 

13 

3 

12 

9 

7 

5 

10 

6 

1 

13 

0 

11 

7 

4 

9 

1 

10 

14 

3 

5 

12 

2 

15 

8 

6 

1 

4 

11 

13 

12 

3 

7 

14 

10 

15 

6 

8 

0 

5 

9 

2 

6 

11 

13 

8 

1 

4 

10 

7 

9 

5 

0 

15 

14 

2 

3 

12 

S8: 

13 

2 

8 

4 

6 

15 

11 

1 

10 

9 

3 

14 

5 

0 

12 

7 

1 

15 

13 

8 

10 

3 

7 

4 

12 

5 

6 

11 

0 

14 

9 

2 

7 

11 

4 

1 

9 

12 

14 

2 

0 

6 

10 

13 

15 

3 

5 

8 

2 

1 

14 

7 

4 

10 

8 

13 

15 

12 

9 

0 

3 

5 

6 

11 


23

Input: 101010

Baris = 10 = 2

Kolom = 0101 = 5

Luaran = 13 = 1101 

• Setiap baris di dalam 

S-box adalah 

permutasi angka-

angka 0, 1, 2,…, 15.

• Mengubah satu bit di 

dalam input S-box 

mengakibatkan 

perubahan paling 

sedikit 2 bit pada 

output S-box

    → prinsip confussion

0      1       2        3       4        5      6        7       8       9       10    11      12     13     14     15 

0

1

2

3

\end{document}
